\section{Overview of the Existing Work}
\textit{TODO (team): summarize existing work and systems.}
this is included


% \section{Overview of Existing Work}
% NOTE:
% We will go over the existing solutions…here.
% For application level: We need to check popular application settled in this scenario….
% Code sharing: Github, gitlab…
% File sharing: Google docs/google driver, Drop Box, Nextcloud
% Directory and ownership sharing: Teams
% P2p sharing: resilio sync, syncthing
% P2p storage: IPFS, Tahoe-LAFS
% For Technical level… we need to go through the existing distributed system and p2p…
% Versioning method: Git, SVN, Perforce..
% P2p protocols
% Sharing: Local Lan, over bluetooth, air drop, …
%

\subsection{Centralized Collaboration Platforms}

% What are centralized collaboration platforms?
% Examples:
% - Google Drive
% - Dropbox
% - OneDrive
% - Nextcloud
%
% Briefly describe:
% - cloud storage
% - file sharing
% - synchronization
% - user management

Centralized collaboration platforms are systems that provide file storage, sharing, and synchronization services through a centrally managed infrastructure.
Users interact with a common server or cloud service that stores project data and coordinates access among collaborators.
These platforms are widely used for collaborative work and document management.

Common functionalities provided by centralized collaboration platforms include cloud-based file storage, folder sharing, cross-device synchronization, and user permission management.
These features enable multiple users to access and modify shared files through a unified platform.

Representative examples include \textit{Google Drive}, \textit{ Dropbox },\textit { Microsoft OneDrive }, and \textit{ Nextcloud }.
Although these systems differ in deployment models and implementation details, they share a common architecture in which a central service manages data storage, synchronization, and user access.

\subsection{Version Control Systems}

% What is version control?
%
% Examples:
% - Git
% - SVN
% - Perforce
%
% Briefly describe:
% - commits
% - revision history
% - rollback
% - collaboration workflow

Version control systems are software tools designed to track changes made to files over time.
They maintain a history of revisions, allowing users to review previous versions, restore earlier states, and coordinate modifications among multiple contributors.
Such systems are commonly used in software development and collaborative content creation.

The main functions of version control systems include change tracking, version history management, conflict handling, and rollback.
By recording how files evolve over time, these systems allow users to inspect previous states, recover from incorrect changes, and coordinate concurrent work within a shared project.

Representative examples include \textit{ Git }, \textit{ Apache Subversion (SVN) }, and \texit{ Perforce }. These systems differ in architecture and workflow, but they all provide functions for recording file changes, preserving historical versions, and supporting collaboration among multiple users.

\subsection{Peer-to-Peer Synchronization Systems}

% What is peer-to-peer synchronization?
%
% Examples:
% - Syncthing
% - Resilio Sync
%
% Briefly describe:
% - peer discovery
% - folder synchronization
% - direct file transfer
% - decentralized architecture


Peer-to-peer synchronization systems are distributed applications that synchronize files directly between participating devices without relying on a central storage server.
Each peer maintains its own copy of the shared data and exchanges updates with other peers in order to keep folders synchronized.

The main functions of peer-to-peer synchronization systems include decentralized data ownership, direct peer-to-peer communication, automatic peer discovery, and distributed replication of shared folders.
By distributing both storage and synchronization responsibilities among participating devices, these systems enable collaborative file sharing without requiring dedicated server infrastructure.

Representative examples include \textit{ Syncthing } and \texit{ Resilio Sync }. Unlike traditional cloud-based solutions, these systems allow participating devices to communicate directly with one another and collectively maintain shared data.


\subsection{Local Sharing Technologies}

% How do devices discover each other and share locally..?
%
% Examples:
% - mDNS
% - Zeroconf
% - Bonjour
% - AirDrop
%
% Briefly describe:
% - service discovery
% - local communication
% - device discovery

Local sharing technologies enable nearby devices to exchange files and other resources directly within a local environment.
Unlike cloud-based collaboration platforms, these technologies typically allow users to share data without first uploading it to a remote server.


The main functions of local sharing technologies include direct device-to-device file transfer, temporary resource sharing, automatic discovery of nearby devices, and simplified local connectivity.
By reducing the need for external infrastructure, these technologies provide a convenient mechanism for local data exchange.

Representative examples include \textit{ AirDrop }, and \textit{ Nearby Share }.
These systems are commonly used for temporary file exchange and temporary collaboration between nearby devices.
\subsection{Related Areas Outside the Scope of This Project}

\subsubsection{Decentralized Storage Systems}

% What is decentralized storage?
%
% Examples:
% - IPFS
% - Tahoe-LAFS
%
% Briefly describe:
% - distributed storage
% - replication
% - content addressing
% - fault tolerance

Decentralized storage systems distribute data across multiple independent nodes rather than storing it on a single server.
These systems aim to improve data availability, fault tolerance, and storage resilience by allowing multiple participants to contribute storage resources to the network.

